\section{Related Literature}
\label{sec:literature}

This paper sits at the intersection of four literatures: the measurement of AI exposure at the occupational level, the post-ChatGPT causal evidence on labor-market outcomes, the task-based theoretical framework for automation and inequality, and the Latin American literature on technology and informal labor markets. I organize the discussion thematically rather than chronologically.

\subsection{Measuring AI exposure}
\label{subsec:lit_exposure}

The canonical exposure measure in the LLM era is \citet{Eloundou2023_gpt_exposure}. Using GPT-4 and a human rubric, the authors rate each O*NET task on whether an LLM (possibly augmented with LLM-powered software) could reduce the time required to perform it by at least 50 percent. Three aggregate occupational scores result: $\alpha$ (direct exposure), $\beta$ (exposure with LLM-powered complementary software), and $\zeta = \alpha + \beta$ (the upper bound). I use $\zeta$ as the primary treatment, with $\alpha$, $\beta$, and alternative measures in robustness. Two alternatives precede Eloundou et al. \citet{Webb2020_ai_exposure} maps AI patents to O*NET tasks, producing an exposure measure that pre-dates generative AI and therefore isolates the pre-LLM automation component. \citet{Felten2023_ai_exposure} construct the AI Occupational Exposure index (AIOE) from expert task-capability assessments in the Electronic Frontier Foundation AI Progress Measurement dataset. These three measures converge on the coarse ranking of occupations (cognitive, language-intensive, clerical/professional work ranks highest) but differ at finer resolutions.

A concern with any predicted exposure measure is whether it reflects actual LLM use. \citet{Handa2025_economic_tasks} address this by analyzing millions of Claude conversations from Anthropic and mapping them to O*NET tasks. The realized-usage ranking correlates strongly with the Eloundou predicted-exposure ranking, which provides empirical content to the ``score-as-dose'' interpretation that I rely on throughout. For Latin America, \citet{Azuara2024_idb_ai} extend the Eloundou measure to Chile, Mexico, and Peru using the BLS SOC-to-ISCO crosswalk, which I build on.

\subsection{Causal effects of LLMs on labor markets}
\label{subsec:lit_causal}

The causal literature is young and growing fast. Firm- and task-level evidence, which comes closest to identifying a clean productivity effect, is generally positive. \citet{Brynjolfsson2023_genai} run a randomized rollout of an LLM-based assistant in a US customer-service call center and document a 14 percent productivity gain concentrated among less-experienced workers. \citet{CuiDemirerJaffe2025_copilot} report 26 percent faster task completion for software developers across three RCTs with nearly 5{,}000 participants. These estimates bound the complementarity channel from above: in a setting where adoption is exogenously delivered, the technology works as advertised.

Aggregate labor-market evidence is more mixed. \citet{Hartley2024_labor_effects} use the US Current Population Survey and a DiD design around the ChatGPT launch to document measurable hours and wage responses in high-exposure occupations. \citet{Chen2025_llm_unemployment} apply a synthetic DiD to US data and find roughly \$89 per week higher earnings in high-exposure occupations, with no detectable effect on unemployment. \citet{Hui2024_online_labor} examine online freelance platforms and show sharp declines (20 to 50 percent) in demand for writing and translation services post-launch, with accompanying earnings drops of 2 to 5 percent. \citet{Liu2025_generate_future} reach similar conclusions in a related online-labor sample. \citet{LiuWangYu2025_labor_demand} use 285 million US job postings and estimate that high-AI-exposure occupations saw a 12 percent decline in posted positions in the first post-launch year, rising to 18 percent by year three.

Against these positive-displacement or positive-earnings estimates, \citet{HumlumVestergaard2025_llm} deliver a striking null. Using Danish administrative data linked to an adoption survey of 100{,}000 workers across eleven exposed occupations, they rule out earnings or hours effects larger than two percent two years post-launch, despite documenting widespread adoption. Their companion paper \citep{HumlumVestergaard2025_pnas} shows that this adoption is itself unequal, skewing male, young, and educated. The tension between positive micro-productivity gains (Brynjolfsson, Cui) and near-null aggregate effects (Humlum-Vestergaard) is the central unresolved question in this literature.

I position this paper as the reference-class challenge to \citet{HumlumVestergaard2025_llm}. Their null may reflect a universal feature of LLM shocks at two-year horizons, or it may reflect specifically Danish labor-market buffers. Latin America offers the opposite structural configuration: high informality, weaker bargaining institutions, and lower but non-trivial adoption. If the null survives in this very different setting, Denmark was not the exception; if it does not, informality is the margin that matters. The paper also serves as external validity for \citet{Hartley2024_labor_effects} and \citet{Chen2025_llm_unemployment}: the harmonized seven-country panel can either confirm their positive estimates in a developing-region setting or document that their magnitudes do not replicate outside the United States.

\subsection{Tasks, automation, and inequality}
\label{subsec:lit_theory}

The theoretical framework follows \citet{AcemogluAutor2011_tasks} and \citet{AcemogluRestrepo2018_race}: technology affects labor demand through two opposing channels. The \emph{displacement effect} substitutes capital (or algorithms) for labor in tasks that the technology can now perform, reducing labor demand in exposed occupations. The \emph{reinstatement effect} creates new tasks or raises the productivity of human labor in tasks complementary to the technology, raising demand in exposed occupations. The net effect on employment and wages is empirical and depends on the relative strength of the two channels. \citet{AcemogluRestrepo2022_tasks_inequality} quantify this framework for US data from 1980--2016 and attribute 50 to 70 percent of the observed rise in wage inequality to task displacement in automating industries, a result this paper's distributional analysis (Paper 2) adapts to the LLM shock.

\citet{AutorLM2003_skill_content} provide the classical routine-biased technical change (RBTC) framework: computers and industrial robots automated \emph{routine} cognitive and manual tasks, compressing wages in the middle of the distribution. LLMs break this pattern. The tasks LLMs most affect - writing, synthesis, analysis, code generation - are \emph{non-routine cognitive} tasks, historically performed by workers at the upper end of the skill distribution. This inverts the classical prediction: high-skill occupations, not middle-skill ones, bear the direct exposure. Whether this translates into displacement of high-skill workers or complementary productivity gains for them is the central empirical question. \citet{Agrawal2019_prediction} formalize the theoretical ambiguity: AI as a prediction technology has inherently mixed effects on labor demand that depend on whether prediction substitutes for or complements human judgment.

The direct empirical precedent for the identification design is \citet{AcemogluAHR2022_ai_jobs}, published in the \emph{Journal of Labor Economics}. They use AI-vacancy data and an occupation-level exposure specification to estimate effects on hiring patterns in US establishments. This paper's design is methodologically similar - occupation-level continuous exposure, occupation and country-time fixed effects - but uses labor force survey microdata rather than vacancy postings, and applies the CGBS 2024 continuous-treatment DiD estimator rather than TWFE.

\subsection{Technology and Latin American labor markets}
\label{subsec:lit_lac}

The distinguishing feature of the Latin American setting is informality. \citet{Ciaschi2025_distributive_ai} estimate the distributive impact of predicted AI exposure using CEDLAS household surveys and document that exposure is concentrated in the upper half of the wage distribution, with informal workers accounting for a disproportionate share of the exposed low-wage population. \citet{EganadelSol2025_ai_lac} examine how generative AI could transform LAC labor markets, and \citet{WorldBank2024_genai_lac} frames the policy discussion. \citet{Georgieff2024_ai_wage_inequality} documents, for OECD countries, that high-exposure occupations exhibit greater wage dispersion - particularly at the top - a motivation for this paper's Paper 2 distributional analysis. \citet{CastellanosRodriguez2024_poverty} and \citet{Brollo2024_social_protection} document the poverty and social-protection context into which the LLM shock lands.

What is missing from this LAC literature is causal identification. Every existing paper stops at descriptive exposure: where the high-exposure workers are, what their demographic characteristics look like, and which segments of the wage distribution they occupy. None estimates the causal effect of LLM exposure on realized labor-market outcomes, because the frontier DiD methods - \citet{CallawayS2021_did}, \citet{CallawayCGBS2024_continuous}, \citet{SunAbraham2021_eventstudy}, \citet{RambachanRoth2023_honest} - have not yet been applied to LLM exposure in a multi-country LAC panel. This paper fills that gap.

\medskip

The paper advances each of these four literatures. In exposure measurement, it validates the Eloundou $\zeta$ score against alternative measures (Webb, Felten, Handa) in a high-informality setting where the score was not constructed. In causal effects, it provides the first non-OECD estimates and the direct comparison with \citet{HumlumVestergaard2025_llm}. In the task-based theoretical framework, it tests whether the displacement-reinstatement decomposition predicts differential effects along the baseline-formality margin that prior work has not examined. In LAC labor economics, it moves the literature from descriptive exposure to causal effects and releases a harmonized seven-country panel for future research.
